\documentclass{article}
\usepackage{graphicx} % Required for inserting images
\usepackage{amsmath,amssymb,amsthm}
\usepackage[a4paper,margin=1in]{geometry}
\title{Proof Of Conjecture 2.3 }
\author{Bishnu Prasad Swain}
\date{June 2026}

\begin{document}
\title{Proof Of Conjecture 2.3 }
\author{Bishnu Prasad Swain}
\date{June 2026}
\maketitle

\noindent\textbf{Notation and known identities}

Let  

\[
\alpha = \theta_2 - \theta_4, \qquad \beta = \theta_3 - \theta_4, \tag{1}
\]  

with \(0<\beta<\alpha<2\pi\) (so \(\theta_2 > \theta_3 > \theta_4\)).  

Define  

\[
p_{24} = \cos\alpha,\quad p_{34} = \cos\beta,\quad p_{23} = \cos(\alpha-\beta). \tag{2}
\]  

From Lemma 2.2 of the paper we have the trigonometric identities  

\[
1-p_{24}-p_{34}+p_{23} = 4\sin\frac{\alpha}{2}\sin\frac{\beta}{2}\cos\frac{\alpha-\beta}{2}, \tag{3}
\]  

\[
1-p_{24} = 2\sin^{2}\frac{\alpha}{2}, \tag{4}
\]  

\[
1-p_{34} = 2\sin^{2}\frac{\beta}{2}. \tag{5}
\]  

Hence the constant  

\[
c = \frac{1-p_{24}-p_{34}+p_{23}}{2\sqrt{(1-p_{24})(1-p_{34})}} = \cos\frac{\alpha-\beta}{2} \; (>0). \tag{6}
\]  

---

\noindent\textbf{Unsquared equation}  

The original geometric condition (equation (15) in the paper) is, with the sign chosen so that both sides are positive,  

\[
\frac{x^{2}-x(p_{24}+p_{34})+p_{23}}{\sqrt{(x^{2}+1-2xp_{24})(x^{2}+1-2xp_{34})}} = c. \tag{7}
\]  

---

\noindent\textbf{Squaring and clearing denominators}  

Squaring (7) and multiplying by the denominator gives  

\[
c^{2}\,(x^{2}+1-2xp_{24})(x^{2}+1-2xp_{34}) = \bigl(x^{2}-x(p_{24}+p_{34})+p_{23}\bigr)^{2}. \tag{8}
\]  

Substitute \(p_{24}=\cos\alpha,\; p_{34}=\cos\beta,\; p_{23}=\cos(\alpha-\beta)\) and define  

\[
P(x)=c^{2}(x^{2}+1-2x\cos\alpha)(x^{2}+1-2x\cos\beta)-\bigl(x^{2}-x(\cos\alpha+\cos\beta)+\cos(\alpha-\beta)\bigr)^{2}. \tag{9}
\]  

Then equation (8) is exactly \(P(x)=0\).

---

\noindent\textbf{Verifying that \(x=\pm1\) are roots}  

For \(x=1\):  

\[
x^{2}+1-2x\cos\alpha = 2-2\cos\alpha = 4\sin^{2}\frac{\alpha}{2}, \tag{10}
\]  

\[
x^{2}+1-2x\cos\beta = 2-2\cos\beta = 4\sin^{2}\frac{\beta}{2}. \tag{11}
\]  

The first term of \(P(1)\) becomes  

\[
c^{2}\cdot 4\sin^{2}\frac{\alpha}{2}\cdot 4\sin^{2}\frac{\beta}{2}=16\,c^{2}\sin^{2}\frac{\alpha}{2}\sin^{2}\frac{\beta}{2}. \tag{12}
\]  

The second term uses  

\[
1-(\cos\alpha+\cos\beta)+\cos(\alpha-\beta)=1-p_{24}-p_{34}+p_{23}=4\sin\frac{\alpha}{2}\sin\frac{\beta}{2}\,c, \tag{13}
\]  

by (3) and (6). Squaring gives  

\[
\bigl(4\sin\frac{\alpha}{2}\sin\frac{\beta}{2}\,c\bigr)^{2}=16\,c^{2}\sin^{2}\frac{\alpha}{2}\sin^{2}\frac{\beta}{2}. \tag{14}
\]  

Thus \(P(1)=0\).  

For \(x=-1\):  

\[
x^{2}+1-2x\cos\alpha = 1+1+2\cos\alpha = 2+2\cos\alpha = 4\cos^{2}\frac{\alpha}{2}, \tag{15}
\]  

\[
x^{2}+1-2x\cos\beta = 1+1+2\cos\beta = 2+2\cos\beta = 4\cos^{2}\frac{\beta}{2}. \tag{16}
\]  

The first term of \(P(-1)\) becomes  

\[
c^{2}\cdot 4\cos^{2}\frac{\alpha}{2}\cdot 4\cos^{2}\frac{\beta}{2}=16\,c^{2}\cos^{2}\frac{\alpha}{2}\cos^{2}\frac{\beta}{2}. \tag{17}
\]  

The second term:  

\[
1+(\cos\alpha+\cos\beta)+\cos(\alpha-\beta)=1+\cos\alpha+\cos\beta+\cos(\alpha-\beta). \tag{18}
\]  

Using \(\cos\alpha=2\cos^{2}\frac{\alpha}{2}-1\) and \(\cos\beta=2\cos^{2}\frac{\beta}{2}-1\), we obtain  

\[
1+\cos\alpha+\cos\beta+\cos(\alpha-\beta)=2\cos^{2}\frac{\alpha}{2}+2\cos^{2}\frac{\beta}{2}+\cos(\alpha-\beta)-1. \tag{19}
\]  

A more compact identity is  

\[
1+\cos\alpha+\cos\beta+\cos(\alpha-\beta)=4\cos\frac{\alpha}{2}\cos\frac{\beta}{2}\cos\frac{\alpha-\beta}{2}, \tag{20}
\]  

which can be verified by standard product‑to‑sum formulas. Since \(\cos\frac{\alpha-\beta}{2}=c\), (20) equals  

\[
4\cos\frac{\alpha}{2}\cos\frac{\beta}{2}\,c. \tag{21}
\]  

Squaring gives  

\[
\bigl(4\cos\frac{\alpha}{2}\cos\frac{\beta}{2}\,c\bigr)^{2}=16\,c^{2}\cos^{2}\frac{\alpha}{2}\cos^{2}\frac{\beta}{2}, \tag{22}
\]  

identical to (17). Hence \(P(-1)=0\).

Therefore \(x=1\) and \(x=-1\) are roots of \(P(x)\), so \((x^{2}-1)\) divides \(P(x)\). Consequently  

\[
P(x)=(x^{2}-1)\,Q(x), \tag{23}
\]  

where \(Q(x)\) is a quadratic polynomial.

---

\noindent\textbf{Explicit form of \(Q(x)\)}  

By expanding \(P(x)\) and factoring out \((x^{2}-1)\) (or by using the trigonometric identities from Lemma 2.3 of the paper), one obtains  

\[
Q(x)=4\sin^{2}\frac{\alpha}{2}\sin^{2}\frac{\beta}{2}\,(x-1)^{2}+4\Bigl(\sin\frac{\alpha}{2}\cos\frac{\beta}{2}-\sin\frac{\beta}{2}\cos\frac{\alpha}{2}\Bigr)^{2}x. \tag{24}
\]  

Observe that  

\[
\sin\frac{\alpha}{2}\cos\frac{\beta}{2}-\sin\frac{\beta}{2}\cos\frac{\alpha}{2}=\sin\frac{\alpha-\beta}{2}\neq0 \quad (\text{since }\alpha\neq\beta). \tag{25}
\]  

Thus the second term in (24) is  

\[
4\sin^{2}\frac{\alpha-\beta}{2}\,x = 4(1-c^{2})\,x, \tag{26}
\]  

which is strictly positive for any \(x>0\). The first term is non‑negative (it is a square times a positive constant). Hence  

\[
Q(x)>0 \quad \text{for all } x>0. \tag{27}
\]  

---

\noindent\textbf{Positive roots of \(P(x)=0\)}  

From (23) and (27), for \(x>0\) we have  

\[
P(x)=0 \;\Longleftrightarrow\; (x^{2}-1)Q(x)=0. \tag{28}
\]  

Since \(Q(x)>0\) for \(x>0\), the only factor that can vanish is \(x^{2}-1\). The positive root of \(x^{2}-1=0\) is \(x=1\). Therefore  

\[
\text{The only positive solution of } P(x)=0 \text{ is } x=1. \tag{29}
\]  

---

\noindent\textbf{Returning to the unsquared equation}  

Any solution of the original unsquared equation (7) automatically satisfies (8) and hence \(P(x)=0\). Moreover, for \(x=1\) the left‑hand side of (7) equals \(c\) (positive), so \(x=1\) is indeed a solution. If there were another positive number \(x_{0}\neq1\) satisfying (7), it would also satisfy \(P(x_{0})=0\), contradicting (29). Thus  

\[
x=1 \text{ is the unique positive solution of equation (7)}. \tag{30}
\]  

---

\noindent\textbf{Conclusion}  

The algebraic factorization \(P(x)=(x^{2}-1)Q(x)\) with \(Q(x)>0\) for \(x>0\) proves that the quartic \(P(x)=0\) has exactly one positive root, namely \(x=1\). This establishes Conjecture 2.3 and consequently Corollary 2.3: the only positive solution of the original geometric condition is \(x=1\).  

The constant \(c\) appears implicitly in the expression of \(Q(x)\) through the relation \(c=\cos\frac{\alpha-\beta}{2}\) (e.g., the second term equals \(4(1-c^{2})x\)), but the crucial property \(Q(x)>0\) for \(x>0\) holds independently.

\end{document}
```

